\section{Modern webfejlesztői csapat felépítése}

Egy modern webfejlesztői csapat több, egymást kiegészítő szerepkörből áll. A csapat célja nem csupán a forráskód megírása, hanem egy teljes webalkalmazás megtervezése, megvalósítása, tesztelése, telepítése és hosszú távú karbantartása. Emiatt a fejlesztési folyamatban technikai, tervezési, minőségbiztosítási és üzemeltetési feladatok is megjelennek \cite{ibmSdlc,atlassianDevOps}.

A folyamat elején általában a product owner vagy projektmenedzser határozza meg az üzleti célokat, a prioritásokat és a fejlesztendő funkciókat \cite{schwaberSutherland2020scrumGuide}. Ő tartja kapcsolatot az ügyféllel vagy a megrendelővel, és segít abban, hogy a csapat ne csak technikailag működő, hanem valódi felhasználói igényeket kielégítő és problémákat megoldó rendszert készítsen.

A felhasználói élmény és a felület megtervezéséért jellemzően a UX/UI designer felel. A UX tervezés a felhasználói útvonalakra, a használhatóságra és az alkalmazás logikus működésére koncentrál, míg a UI tervezés a vizuális megjelenést, az elrendezést, a színeket, a tipográfiát és az interakciós elemeket alakítja ki \cite{nielsenNorman1998ux}. Egy jól megtervezett felület csökkenti a fejlesztési bizonytalanságot, mert a fejlesztők pontosabb képet kapnak arról, mit kell megvalósítaniuk.

A frontend fejlesztő a böngészőben futó kliensoldali részt készíti el. Ide tartozik a felhasználói felület HTML, CSS és JavaScript alapú megvalósítása, valamint modern keretrendszerek, például React, Angular vagy Vue használata \cite{mdnWebStandardsModel,mdnClientFrameworks}. A frontend fejlesztő feladata, hogy a felület reszponzív, gyors, hozzáférhető és felhasználóbarát legyen, miközben megfelelően kommunikál a backend által biztosított API-kkal \cite{mdnResponsiveDesign,w3cWaiAccessibilityIntro}.

A backend fejlesztő a szerveroldali logikáért felel. Feladata az üzleti logika megvalósítása, az adatbázis-kezelés, az autentikáció, az autorizáció, valamint az API-k kialakítása \cite{mdnServerSideIntro,mdnClientServerOverview}. A backend biztosítja, hogy a rendszer megbízhatóan kezelje az adatokat, megfelelő válaszokat adjon a kliensoldali kérésekre, és biztonságosan működjön több felhasználó vagy nagyobb terhelés mellett is (skálázhatóság).

Nagyobb csapatokban megjelenhet a full-stack fejlesztő szerepköre is, aki frontend és backend feladatokat egyaránt képes ellátni \cite{awsFullStackDevelopment}. Ez különösen kisebb projektekben hasznos, ahol nincs külön ember minden technológiai rétegre. Ugyanakkor nagyobb rendszereknél a specializáció általában hatékonyabb, mert a kliensoldali és szerveroldali fejlesztés is külön mély szakértelmet igényel.

A minőségbiztosításért a QA engineer vagy tesztelő felel. Ő manuális és automatizált tesztekkel ellenőrzi, hogy az alkalmazás megfelel-e a követelményeknek, és nem tartalmaz-e kritikus hibákat. A tesztelés nem csupán a fejlesztés végén fontos, hanem a teljes folyamat során, mivel a korán felismert hibák javítása általában olcsóbb és egyszerűbb \cite{ibmShiftLeftTesting}. Különösen fontosak a regressziós hibák kiszűrésére szolgáló tesztek, mivel egy új fejlesztés vagy hibajavítás könnyen elronthat egy korábban már működő funkciót \cite{istqbGlossaryRegressionTesting}. Az ismételhető, automatizált tesztek ezt a kockázatot csökkentik \cite{istqbGlossaryTestAutomation}.

A DevOps engineer vagy üzemeltetési szakember a fejlesztés és az infrastruktúra közötti kapcsolatot biztosítja. Feladata lehet a CI/CD folyamatok kialakítása, a konténerizáció, a felhőalapú infrastruktúra kezelése, a monitoring, a naplózás és az automatikus telepítések beállítása \cite{atlassianDevOps,awsCicd}. A DevOps szemlélet célja, hogy a fejlesztés, tesztelés és telepítés gyorsabb, megbízhatóbb és automatizáltabb legyen.

A modern webfejlesztői csapat működése erősen együttműködésre épül \cite{agileManifesto2001}. A szerepkörök nem teljesen elszigeteltek: a frontend fejlesztőnek értenie kell az API-k működését, a backend fejlesztőnek figyelembe kell vennie a felhasználói igényeket, a designernek ismernie kell a technikai korlátokat, a tesztelőnek pedig már a fejlesztés korai szakaszában részt kell vennie a követelmények pontosításában. Ezért a hatékony webfejlesztés nemcsak technológiai, hanem kommunikációs és szervezési feladat is.
